\subsubsection{Classical Ising Model}\label{sec:classical-ising-model}

The \definition{Ising Model} is a \textbf{mathematical model} originally introduced in statistical mechanics to describe \textbf{ferromagnetism}, that is, how microscopic magnetic particles interact with one another.

\highspace
The central idea is extremely simple:
\begin{itemize}
    \item We have a collection of particles.
    \item Each particle behaves like a tiny magnet.
    \item Each magnet can point either up ($+1$) or down ($-1$).
\end{itemize}
The value describing \textbf{where the tiny magnet points} is the \definition{Spin} variable, denoted as $s_i$ for the $i$-th particle. Unlike QUBO variables, which are binary:
\begin{equation*}
    x_i \in \{0,1\}
\end{equation*}
The \textbf{Ising model} uses \textbf{spin variables} that can take values of either $+1$ or $-1$:
\begin{equation*}
    s_i \in \{-1, +1\}
\end{equation*}

\highspace
\begin{flushleft}
    \textcolor{Green3}{\faIcon{book} \textbf{The Ising Energy Function}}
\end{flushleft}
The energy of the system is given by the \definition{Ising Energy Function}:
\begin{equation}\label{eq:ising-energy-function}
    H_{\text{Ising}} = \sum_i h_i s_i + \sum_{i>j} J_{ij} s_i s_j
\end{equation}
At first glance it looks complicated, but it is actually composed of two very simple parts.
\begin{itemize}
    \item \important{Local Fields}:
    \begin{equation*}
        \sum_i h_i s_i
    \end{equation*}
    The coefficient $h_i$ represents an external magnetic field acting on spin $i$. We can think of it as a preference. This \textbf{coefficient encourages the spin to align in a certain direction}:
    \begin{itemize}
        \item If \hl{$h_i > 0$}, the field encourages $s_i$ to be $-1$ (spin \textbf{down}).
        \item If \hl{$h_i < 0$}, the field encourages $s_i$ to be $+1$ (spin \textbf{up}).
    \end{itemize}
    \textcolor{Green3}{\faIcon{question-circle} \textbf{Why does a positive $h_i$ encourage $s_i$ to be $-1$? Isn't that counterintuitive?}} Seemingly, it might be counterintuitive at first, but remember that the \textbf{energy function is something we want to minimize}. If $h_i > 0$, then having $s_i = -1$ will contribute a negative term to the energy, thus lowering it. Conversely, if $s_i = +1$, it would contribute a positive term, increasing the energy. Therefore, a positive $h_i$ encourages $s_i$ to be $-1$ to minimize the energy.

    Suppose $h_i = 5$, then the contribution to the energy from this term would be:
    \begin{itemize}
        \item If $s_i = +1 \to 5 \cdot (+1) = +5$ (increases energy)
        \item If $s_i = -1 \to 5 \cdot (-1) = -5$ (decreases energy)
    \end{itemize}
    Thus, the system will \emph{prefer} $s_i = -1$ to minimize the energy when $h_i$ is positive.

    \textcolor{Green3}{\faIcon{question-circle} \textbf{Why is it called ``\emph{local fields}''?}} The word \textbf{local} is there because this term acts on \textbf{one spint at a time}. Each $h_i$ only affects the energy contribution of the corresponding spin $s_i$, without directly involving any other spins. In contrast, the second term involves interactions between pairs of spins, which is why we call it \textbf{spin interactions}.


    \item \important{Spin Interactions}:
    \begin{equation*}
        \sum_{i>j} J_{ij} s_i s_j
    \end{equation*}
    This term describes \textbf{how pairs of spins influence each other}. The coefficient $J_{ij}$ is called the \textbf{coupling strength} between spins $i$ and $j$.
    \begin{itemize}
        \item[\textcolor{Green3}{\faIcon{question-circle}}] \textcolor{Green3}{\textbf{Why $i>j$ and not $i \neq j$?}} The notation $i>j$ is just a convention to avoid double-counting the interactions. Suppose we have 3 spins: $s_1$, $s_2$, and $s_3$; if we use $i \neq j$, we would get:
        \begin{equation*}
            J_{12} s_1 s_2 + \underbrace{J_{21} s_2 s_1}_{\text{duplicate}} + J_{13} s_1 s_3 + \underbrace{J_{31} s_3 s_1}_{\text{duplicate}} + J_{23} s_2 s_3 + \underbrace{J_{32} s_3 s_2}_{\text{duplicate}}
        \end{equation*}
        But since $s_i s_j$ is the same as $s_j s_i$, we would be \textbf{counting each interaction twice}.

        \item[\textcolor{Green3}{\faIcon{question-circle}}] \textcolor{Green3}{\textbf{What does $J_{ij}$ represent?}} The coefficient $J_{ij}$ represents the strength of the interaction between spins $i$ and $j$.

        Remember $s_{i}, s_{j} \in \{-1, +1\}$, so the product $s_i s_j$ can only be:
        \begin{equation*}
            s_i s_j = \begin{cases}
                +1 & \text{same direction} \\
                -1 & \text{opposite directions}
            \end{cases}
        \end{equation*}
        Because:
        \begin{center}
            \begin{tabular}{@{} c c c @{}}
                \toprule
                $s_i$ & $s_j$ & $s_i s_j$ \\
                \midrule
                $+1$ & $+1$ & $+1$ \\
                $+1$ & $-1$ & $-1$ \\
                $-1$ & $+1$ & $-1$ \\
                $-1$ & $-1$ & $+1$ \\
                \bottomrule
            \end{tabular}
        \end{center}
        Thus, the value of $J_{ij}$ determines \hl{whether the spins prefer to align in the same direction or in opposite directions}. Remember that we want to \textbf{minimize the energy}, therefore:
        \begin{itemize}
            \item \hl{$J_{ij} > 0$}: if the spins are aligned in the same direction ($s_i s_j = +1$), the energy contribution is $+J_{ij}$, which increases the energy. If they are in opposite directions ($s_i s_j = -1$), the contribution is $-J_{ij}$, which decreases the energy. Therefore, the system prefers $s_i s_j = -1$, meaning the \textbf{spins prefer to align in opposite directions} (one up and one down).
            \item \hl{$J_{ij} < 0$}: if the spins are aligned in the same direction ($s_i s_j = +1$), the energy contribution is $-J_{ij}$ (since $J_{ij}$ is negative, this actually decreases the energy). If they are in opposite directions ($s_i s_j = -1$), the contribution is $+J_{ij}$ (which increases the energy). Therefore, the system prefers $s_i s_j = +1$, meaning the \textbf{spins prefer to align in the same direction} (both up or both down).
        \end{itemize}
    \end{itemize}

    In contrast to the \emph{local fields}, which act on individual spins, the \textbf{spin interactions} describe how pairs of spins affect each other.
\end{itemize}
\textcolor{Green3}{\faIcon{question-circle} \textbf{Where do $h_i$ and $J_{ij}$ come from?}} In the original physical context of the Ising model, $h_i$ represents an external magnetic field applied to the system, while $J_{ij}$ represents the intrinsic interaction strength between neighboring spins. Therefore, they are parameters that physicists would determine based on the physical properties of the material being studied.

\highspace
However, in quantum optimization, we can treat $h_i$ and $J_{ij}$ as \textbf{tunable parameters} that we can set to encode specific optimization problems. By choosing appropriate values for $h_i$ and $J_{ij}$, we can represent a wide variety of combinatorial optimization problems within the framework of the Ising model.

\highspace
\textcolor{Green3}{\faIcon{question-circle} \textbf{If $h_i$ and $J_{ij}$ are tunable parameters, are the spins $s_i$ also tunable?}} No, the spin variables $s_i$ are \textbf{not tunable parameters}; they are the \textbf{variables we want to optimize over}. The goal is to find the configuration of spins (i.e., the values of $s_i$) that minimizes the energy function $H_{\text{Ising}}$. In other words, we choose $h_i$ and $J_{ij}$ so that the energy function represents our optimization problem, and then we \hl{search for the spin configuration that yields the lowest energy, which corresponds to the optimal solution of the problem we are trying to solve.}

\highspace
In other words, in optimization problems, the coefficients $h_i$ and $J_{ij}$ define the energy landscape, while the spin variables $s_i$ represent the candidate solution. Solving the problem means finding the spin configuration that minimizes the Ising Hamiltonian.

\highspace
\begin{flushleft}
    \textcolor{Green3}{\faIcon{question-circle} \textbf{Wait, the Ising model equation looks a lot like the QUBO energy function. Are they related?}}
\end{flushleft}
The QUBO energy function is given by (\autopageref{eq:qubo-formulation}):
\begin{equation*}
    E(x) = \sum_i a_i x_i + \sum_{i>j} q_{ij} x_i x_j
\end{equation*}
Similarly, the Ising energy function is given by:
\begin{equation*}
    H_{\text{Ising}} = \sum_i h_i s_i + \sum_{i>j} J_{ij} s_i s_j
\end{equation*}
Notice that they have the same structure: a linear term (the first sum) and a quadratic term (the second sum).
\begin{center}
    \begin{tabular}{@{} c c @{}}
        \toprule
        \textbf{QUBO} & \textbf{Ising} \\
        \midrule
        $x_i \in \{0,1\}$ & $s_i \in \{-1,+1\}$ \\
        $a_i$ & $h_i$ \\
        $q_{ij}$ & $J_{ij}$ \\
        \bottomrule
    \end{tabular}
\end{center}

\noindent
This similarity is really important, because the \hl{Ising energy function is essentially another way of writing an optimization problem, just like the QUBO energy function.} The only \textbf{major difference is the domain of the variables}: QUBO uses binary variables ($0$ and $1$), while the Ising model uses spin variables ($-1$ and $+1$).